\section*{NeurIPS Paper Checklist}

\begin{enumerate}

\item {\bf Claims}
    \item[] Question: Do the main claims made in the abstract and introduction accurately reflect the paper's contributions and scope?
    \item[] Answer: \answerYes{}
    \item[] Justification: The abstract and Section~\ref{sec:intro} state the geometric lower bound (diagonal cost), the regime-based router, and the empirical contract; Sections~\ref{sec:setup}--\ref{sec:empirical} establish each. Falsifiable predictions are stated explicitly in Section~1 under ``Scope and Falsifiable Predictions.''

\item {\bf Limitations}
    \item[] Question: Does the paper discuss the limitations of the work performed by the authors?
    \item[] Answer: \answerYes{}
    \item[] Justification: Section~\ref{sec:limitations} documents encoder blindness, implicit-$k$ extraction, scope to explicit constraints, and the 11\% pivot regime as known failure modes; Appendix~D.1 reports robustness stress tests under paraphrasing, reordering, and softening.

\item {\bf Theory assumptions and proofs}
    \item[] Question: For each theoretical result, does the paper provide the full set of assumptions and a complete (and correct) proof?
    \item[] Answer: \answerYes{}
    \item[] Justification: All theorems state assumptions explicitly (Hilbert space, finite $k$, displacement contract). Full proofs appear in Appendix~\ref{app:proofs}; proof sketches appear inline in the main text.

\item {\bf Experimental result reproducibility}
    \item[] Question: Does the paper fully disclose all the information needed to reproduce the main experimental results of the paper to the extent that it affects the main claims and/or conclusions of the paper (regardless of whether the code and data are provided or not)?
    \item[] Answer: \answerYes{}
    \item[] Justification: Four deterministic verifiers (code AST, JSON schema, optimization gradient, signal DSL) replace LLM-as-judge. Pre-computed results in supplementary JSON files enable direct inspection without re-running expensive experiments. Verification harnesses listed in Appendix~\ref{app:harness_suite}.

\item {\bf Open access to data and code}
    \item[] Question: Does the paper provide open access to the data and code, with sufficient instructions to faithfully reproduce the main experimental results, as described in supplemental material?
    \item[] Answer: \answerYes{}
    \item[] Justification: Anonymized repository contains all verification harnesses, embeddings, and reproducibility scripts. Reproducibility guide maps every quantitative claim to verification code (CLAIM\_AUDIT.md).

\item {\bf Experimental setting/details}
    \item[] Question: Does the paper specify all the training and test details (e.g., data splits, hyperparameters, how they were chosen, type of optimizer) necessary to understand the results?
    \item[] Answer: \answerYes{}
    \item[] Justification: Appendix~\ref{app:experiments} documents data splits, models tested, encoder choices, and routing thresholds. Table~\ref{tab:lipschitz_calibration} reports Lipschitz calibration per model. Routing thresholds are pre-registered in Algorithm~\ref{alg:routing}.

\item {\bf Experiment statistical significance}
    \item[] Question: Does the paper report error bars suitably and correctly defined or other appropriate information about the statistical significance of the experiments?
    \item[] Answer: \answerYes{}
    \item[] Justification: Lipschitz calibration reports mean $\pm$ standard deviation per model (Table~\ref{tab:lipschitz_calibration}); rank correlations report Spearman $r_s$, Kendall $\tau$, sample size, and p-values (e.g., $r_s=-0.942$, 48 points, $p=2.1\times10^{-23}$, Section~\ref{sec:empirical}).

\item {\bf Experiments compute resources}
    \item[] Question: For each experiment, does the paper provide sufficient information on the computer resources (type of compute workers, memory, time of execution) needed to reproduce the experiments?
    \item[] Answer: \answerYes{} % [VERIFY: confirm and add concrete totals before camera-ready]
    \item[] Justification: Inference-time experiments use commercial LLM APIs (GPT-4o, Claude family, Gemini, DeepSeek) and local inference for open-weight models (Qwen-2.5-Coder, TinyLlama-1.1B, CodeLlama-7B) on a single workstation. Embedding computation is CPU-only ($<$10\,ms per regime-index call). Total API spend on the order of \$[XXX] over the experimental period; per-experiment costs noted in Appendix~\ref{app:experiments}. % [VERIFY $XXX]

\item {\bf Code of ethics}
    \item[] Question: Does the research conducted in the paper conform, in every respect, with the NeurIPS Code of Ethics?
    \item[] Answer: \answerYes{}
    \item[] Justification: The research conforms with the NeurIPS Code of Ethics. No human subjects, no scraped or sensitive data, no released models requiring additional safeguards.

\item {\bf Broader impacts}
    \item[] Question: Does the paper discuss both potential positive societal impacts and negative societal impacts of the work performed?
    \item[] Answer: \answerYes{}
    \item[] Justification: \textbf{Positive:} Making geometric limits visible is prerequisite to respecting them; practitioners cannot avoid cliffs they cannot detect. The pre-generation routing rule lets aligned systems decline impossible requests rather than degrade silently. \textbf{Risk:} Same principles could optimize harmful constraint satisfaction (e.g., compound adversarial prompts). \textbf{Mitigation:} The bound is symmetric---adversaries face the same limits. $\hat{\rho}$ detection works bidirectionally: compound adversarial prompts are \emph{more} visible (same cliffs). Fail-safe property: when constraints conflict, the system abstains rather than persuades. \textbf{Commercial relevance:} Advertising creates constraint coupling (user intent vs.\ promotion vs.\ transparency); the diagonal cost predicts infeasibility, informing context-integrity policy.

\item {\bf Safeguards}
    \item[] Question: Does the paper describe safeguards that have been put in place for responsible release of data or models that have a high risk for misuse (e.g., pre-trained language models, image generators, or scraped datasets)?
    \item[] Answer: \answerNA{}
    \item[] Justification: The paper releases no models, datasets, or assets with misuse risk. The contribution is a geometric lower bound and a routing algorithm using existing publicly available embedding models.

\item {\bf Licenses for existing assets}
    \item[] Question: Are the creators or original owners of assets (e.g., code, data, models), used in the paper, properly credited and are the license and terms of use explicitly mentioned and properly respected?
    \item[] Answer: \answerYes{} % [VERIFY: license claims]
    \item[] Justification: All evaluated LLMs are accessed via published APIs (commercial models) or open-weight releases under their respective licenses (Llama 3.1 community license, Qwen Apache 2.0, TinyLlama Apache 2.0, CodeLlama community license). Sentence-Transformers encoders (MiniLM-L6-v2, mpnet-base-v2, multilingual-MiniLM) are Apache 2.0. Each is cited in the relevant section.

\item {\bf New assets}
    \item[] Question: Are new assets introduced in the paper well documented and is the documentation provided alongside the assets?
    \item[] Answer: \answerNA{}
    \item[] Justification: No new datasets or pre-trained models are released. Verification harnesses are released as code with documentation in the anonymized repository.

\item {\bf Crowdsourcing and research with human subjects}
    \item[] Question: For crowdsourcing experiments and research with human subjects, does the paper include the full text of instructions given to participants and screenshots, if applicable, as well as details about compensation (if any)?
    \item[] Answer: \answerNA{}
    \item[] Justification: The research does not involve crowdsourcing or human subjects.

\item {\bf Institutional review board (IRB) approvals or equivalent for research with human subjects}
    \item[] Question: Does the paper describe potential risks incurred by study participants, whether such risks were disclosed to the subjects, and whether Institutional Review Board (IRB) approvals (or an equivalent approval/review based on the requirements of your country or institution) were obtained?
    \item[] Answer: \answerNA{}
    \item[] Justification: No human subjects research.

\item {\bf Declaration of LLM usage}
    \item[] Question: Does the paper describe the usage of LLMs if it is an important, original, or non-standard component of the core methods in this research?
    \item[] Answer: \answerYes{}
    \item[] Justification: LLMs are the \emph{object of study} in this paper, not a writing aid for the methodology. The empirical work analyzes multi-constraint generation behavior across GPT-4o, Claude (haiku-3 through opus-4.5), Gemini-2.5 Flash and Pro, DeepSeek-v3 and -R1, Llama-3.1-70B, Qwen-2.5-Coder, TinyLlama-1.1B, and CodeLlama-7B. Specific models, versions, and access methods are documented in Section~\ref{sec:empirical} and Appendix~\ref{app:experiments}. LLM-assisted writing was used for prose editing during manuscript preparation; this does not affect the core methodology, scientific rigor, or originality of the research.

\end{enumerate}
